% !TeX root = main.tex
% PAGINE ORIGINALI 481-490 — DA COMPILARE
% Incollare qui i quiz del blocco e alla fine le rispettive\registraSoluzione{numero}{lettera}.
% Questo file vuoto non aggiunge pagine al PDF.

% !TeX root = main.tex
% Pagine 481-490: 5 quesiti della 28a simulazione e 20 quesiti della 29a simulazione.

% Pag. originale 481/521
\iffalse Pag. originale 481/521 \fi
\begin{quiz}
$\displaystyle\lim_{x\to+\infty}x\left(\cos\frac1x-e^{4/x}\right)=$
\begin{enumerate}
\item $-4$
\item $1$
\item $\nexists$
\item $\dfrac12$
\item $-\infty$
\end{enumerate}
\end{quiz}
% SOLUZIONE: A
% SEZIONE_AM1: 3.4
% SPIEGAZIONE: Ponendo $t=1/x$, si ha $\cos t-e^{4t}=-4t+O(t^2)$, dunque il limite vale $-4$.

% Pag. originale 481/521
\iffalse Pag. originale 481/521 \fi
\begin{quiz}
L'equazione differenziale $y'=e^{x+y}$:
\begin{enumerate}
\item \`e lineare
\item \`e a variabili separabili
\item ammette $y=e^x$ come soluzione particolare
\item la soluzione che soddisfa il problema di Cauchy $y(0)=1$ \`e $y=e^x$
\item nessuna delle precedenti
\end{enumerate}
\end{quiz}
% SOLUZIONE: B
% SEZIONE_AM1: 6.1
% SPIEGAZIONE: $y\prime=e^x e^y$ separa le variabili: $e^{-y}\,dy=e^x\,dx$. Non \`e lineare in $y$.

% Pag. originale 481/521
\iffalse Pag. originale 481/521 \fi
\begin{quiz}
$\displaystyle\lim_{n\to+\infty}(3+n)^{5/\log n}=$
\begin{enumerate}
\item $5$
\item $1$
\item $+\infty$
\item $0$
\item $e^5$
\end{enumerate}
\end{quiz}
% SOLUZIONE: E
% SEZIONE_AM1: 3.4
% SPIEGAZIONE: Il logaritmo della potenza \`e $5\log(n+3)/\log n\to5$; per continuit\`a della funzione esponenziale il limite \`e $e^5$.

% Pag. originale 481/521
\iffalse Pag. originale 481/521 \fi
\begin{quiz}
Senza il calcolo della derivata prima, si pu\`o affermare che la funzione $f:[0,2\pi]\to\mathbb R$, $f(x)=e^{\cos x}\sin x$:
\begin{enumerate}
\item ammette almeno due punti critici in $(0,2\pi)$
\item ha un solo punto critico in $(0,2\pi)$
\item non ha punti critici in $(0,\pi)$
\item non ha punti critici in $(\pi,2\pi)$
\item $f'(x)=0$ ha una sola soluzione con $x\in(0,2\pi)$
\end{enumerate}
\end{quiz}
% SOLUZIONE: A
% SEZIONE_AM1: 4.3
% SPIEGAZIONE: $f(0)=f(\pi)=f(2\pi)=0$. Per Rolle, esiste un punto stazionario in $(0,\pi)$ e uno in $(\pi,2\pi)$.

% Pag. originale 481/521
\iffalse Pag. originale 481/521 \fi
\begin{quiz}
$\displaystyle\lim_{x\to4/5}\frac{1-\cos(5x-4)}{(x-\frac45)^2}=$
\begin{enumerate}
\item $\dfrac12$
\item $\dfrac{25}{2}$
\item $\dfrac52$
\item $1$
\item $\dfrac1{5\cdot2}$
\end{enumerate}
\end{quiz}
% SOLUZIONE: B
% SEZIONE_AM1: 3.4
% SPIEGAZIONE: Con $t=5x-4=5(x-4/5)$, il limite \`e $25\lim_{t\to0}(1-\cos t)/t^2=25/2$.

% Pag. originale 486/521
\iffalse Pag. originale 486/521 \fi
\begin{quiz}
Siano $A$ e $B$ due insiemi, con $A$ e $B\subseteq\mathbb R$ e tali che $\sup A=\sup B=3$. Allora
\begin{enumerate}
\item $\max A=\max B$
\item se $\max A=3$ allora $\max B=3$
\item $A\equiv B$
\item $\inf A=\inf B$
\item $B$ pu\`o non ammettere massimo
\end{enumerate}
\end{quiz}
% SOLUZIONE: E
% SEZIONE_AM1: 1.5
% SPIEGAZIONE: L\'estremo superiore pu\`o non appartenere a $B$, per esempio $B=(0,3)$, che non ha massimo.

% Pag. originale 486/521
\iffalse Pag. originale 486/521 \fi
\begin{quiz}
Sia $A\subseteq\mathbb R$ e $\inf A=5$
\begin{enumerate}
\item $5$ \`e uno degli elementi di $A$
\item non \`e detto che $5$ sia il minimo di $A$
\item $5$ \`e il pi\`u piccolo dei minoranti di $A$
\item $5\in A$
\item esistono elementi di $A$ minori di $5$
\end{enumerate}
\end{quiz}
% SOLUZIONE: B
% SEZIONE_AM1: 1.5
% SPIEGAZIONE: L\'infimo \`e il maggiore dei minoranti, ma non deve appartenere all\'insieme. Per esempio $A=(5,6)$.

% Pag. originale 486/521
\iffalse Pag. originale 486/521 \fi
\begin{quiz}
Sia $f:\mathbb R\to\mathbb R$ una funzione continua tale che $f(x)=-f(-x)$. Quale delle seguenti affermazioni \`e sempre vera?
\begin{enumerate}
\item $\displaystyle\int_{-1}^{1}f(x)\,dx=2\int_0^1f^2(x)\,dx$
\item $\displaystyle\int_0^1f(x)\,dx=0$
\item $\displaystyle\int_{-1}^1f(x)\,dx=0$
\item $\displaystyle\int_{-1}^1f(x)\,dx=2\int_0^1f(x)\,dx$
\item $\displaystyle\int_{-1}^0f(x)\,dx=\int_0^1f(x)\,dx$
\end{enumerate}
\end{quiz}
% SOLUZIONE: C
% SEZIONE_AM1: 5.2
% SPIEGAZIONE: La funzione \`e dispari; l\'integrale su un intervallo simmetrico rispetto all\'origine vale zero.

% Pag. originale 486/521
\iffalse Pag. originale 486/521 \fi
\begin{quiz}
Sia $y(x)$ la soluzione del problema di Cauchy $\begin{cases}y\prime\prime+4y=0,\\y(0)=1,\\y\prime(0)=0.\end{cases}$ Allora
\begin{enumerate}
\item $\displaystyle\lim_{x\to+\infty}y(x)=0$
\item $\displaystyle\lim_{x\to+\infty}y(x)=-\infty$
\item $\displaystyle\lim_{x\to+\infty}y(x)=\nexists$
\item $\displaystyle\lim_{x\to+\infty}y(x)=1$
\item $\displaystyle\lim_{x\to+\infty}y(x)=\pi$
\end{enumerate}
\end{quiz}
% SOLUZIONE: C
% SEZIONE_AM1: 6.1
% SPIEGAZIONE: La soluzione del problema \`e $y(x)=\cos(2x)$, che oscilla e non ha limite per $x\to+\infty$.

% Pag. originale 487/521
\iffalse Pag. originale 487/521 \fi
\begin{quiz}
Se $g(x)=f(x^2)$, allora il polinomio di Taylor di grado $2$ della funzione $g$ con centro in $x=0$ \`e:
\begin{enumerate}
\item $f(0)+f\prime(0)x+f\prime\prime(0)x^2$
\item $2f\prime(0)x+f\prime\prime(0)x^2$
\item $f(0)+f\prime(0)x+\dfrac12f\prime\prime(0)x^2$
\item $f(0)+\dfrac1{2!}f\prime(0)x^2$
\item $f(0)+f\prime(0)x^2$
\end{enumerate}
\end{quiz}
% SOLUZIONE: E
% SEZIONE_AM1: 3.3
% SPIEGAZIONE: $f(x^2)=f(0)+f\prime(0)x^2+o(x^2)$, assumendo la regolarit\`a necessaria. La risposta D differisce per il fattore $1/2$.

% Pag. originale 487/521
\iffalse Pag. originale 487/521 \fi
\begin{quiz}
Sia $f:\mathbb R\to\mathbb R$ una funzione di classe $C^{(n)}(\mathbb R)$. Quale delle seguenti affermazioni \`e SEMPRE vera?
\begin{enumerate}
\item Se $f(1)=1$ e $f(-1)=-1$ allora $f^2$ ha un minimo assoluto.
\item Se $f$ ha un punto di minimo allora $f^2$ ha un punto di minimo
\item Se $f$ \`e crescente allora $f^2$ \`e crescente
\item Se $f$ \`e convessa allora $f^2$ \`e convessa.
\item Se $f(1)=1$ e $f(-1)=-1$ allora $f^3$ ha un minimo assoluto.
\end{enumerate}
\end{quiz}
% SOLUZIONE: A
% SEZIONE_AM1: 4.3
% SPIEGAZIONE: Per il teorema degli zeri, $f$ si annulla in un punto di $(-1,1)$. Quindi $f^2\ge0$ raggiunge il minimo assoluto $0$.

% Pag. originale 487/521
\iffalse Pag. originale 487/521 \fi
\begin{quiz}
$\displaystyle\int_0^1xf(2x^2+3)\,dx=$
\begin{enumerate}
\item $\displaystyle\frac14\int_3^5f(t)\sqrt{\frac{t-3}{2}}\,dt$
\item $\displaystyle4\int_3^5f(t)\sqrt{\frac{t-3}{2}}\,dt$
\item $\displaystyle\frac14\int_3^5f(t)\,dt$
\item $\displaystyle4\int_3^5f(t)\,dt$
\item $\displaystyle4\int_0^1f(t+3)\,dt$
\end{enumerate}
\end{quiz}
% SOLUZIONE: C
% SEZIONE_AM1: 5.4
% SPIEGAZIONE: Con $t=2x^2+3$, $dt=4x\,dx$; gli estremi diventano $3$ e $5$.

% Pag. originale 487/521
\iffalse Pag. originale 487/521 \fi
\begin{quiz}
Lo sviluppo di Mac Laurin, di ordine $6$, della funzione $f(x)=\log\bigl(2+\sin(x^6)\bigr)$ \`e
\begin{enumerate}
\item $\log2+x^6+o(x^6)$
\item $1+\dfrac{x^6}{2}+o(x^6)$
\item $2+x^6+o(x^6)$
\item $\log2+\dfrac{x^6}{2}+o(x^6)$
\item $\log2+\dfrac{x^6}{6!}+o(x^6)$
\end{enumerate}
\end{quiz}
% SOLUZIONE: D
% SEZIONE_AM1: 3.3
% SPIEGAZIONE: $\sin(x^6)=x^6+o(x^6)$ e $\log(2+u)=\log2+u/2+o(u)$.

% Pag. originale 487/521
\iffalse Pag. originale 487/521 \fi
\begin{quiz}
$\displaystyle\lim_{x\to+\infty}\frac3x\log(e^{2x}+x)=$
\begin{enumerate}
\item $6$
\item $0$
\item $2$
\item $+\infty$
\item $\infty$
\end{enumerate}
\end{quiz}
% SOLUZIONE: A
% SEZIONE_AM1: 3.4
% SPIEGAZIONE: $\log(e^{2x}+x)=2x+\log(1+xe^{-2x})=2x+o(1)$; quindi il limite vale $6$.

% Pag. originale 488/521
\iffalse Pag. originale 488/521 \fi
\begin{quiz}
Sia data la funzione $f(x)=\sinh x$. Allora:
\begin{enumerate}
\item $f^{(n)}(0)=1$ solo se $n=4k+1$, $k\in\mathbb N$
\item $f^{(n)}(\pi/2)=1$ se $n\ge2$, dispari
\item $f^{(n)}(0)=-1$ se $n\ge2$, pari
\item $f^{(n)}(0)=1$ se $n=2k+1$, $k\in\mathbb N$
\item $f^{(n)}(\pi/2)=1$ se $k=4k$, $k\in\mathbb N$
\end{enumerate}
\end{quiz}
% SOLUZIONE: D
% SEZIONE_AM1: 3.1
% SPIEGAZIONE: Le derivate di $\sinh x$ alternano $\cosh x$ e $\sinh x$; quelle dispari in zero valgono $1$, quelle pari valgono $0$.

% Pag. originale 488/521
\iffalse Pag. originale 488/521 \fi
\begin{quiz}
Sia $F$ una primitiva della funzione $f:\mathbb R\to\mathbb R$, continua. Sappiamo che $F$ ha un massimo in $x_0$, allora necessariamente
\begin{enumerate}
\item $f$ si annulla in $x_0$
\item $f$ \`e derivabile in $x_0$
\item $f$ ha un minimo in $x_0$
\item $F$ si annulla in $x_0$
\item $f\ge0$, $\forall x\in\mathbb R$
\end{enumerate}
\end{quiz}
% SOLUZIONE: A
% SEZIONE_AM1: 5.1
% SPIEGAZIONE: Essendo $F\prime=f$ e $F$ derivabile, per Fermat $F\prime(x_0)=f(x_0)=0$.

% Pag. originale 488/521
\iffalse Pag. originale 488/521 \fi
\begin{quiz}
$\displaystyle\lim_{x\to4/5}\frac{1-\cos(5x-4)}{(x-\frac45)^2}=$
\begin{enumerate}
\item $\dfrac{25}{2}$
\item $\dfrac{25}{4}$
\item $25$
\item $\dfrac45$
\item $1$
\end{enumerate}
\end{quiz}
% SOLUZIONE: A
% SEZIONE_AM1: 3.4
% SPIEGAZIONE: Come nel quesito 20 della 28a simulazione, $1-\cos t\sim t^2/2$ con $t=5(x-4/5)$, dunque $25/2$.

% Pag. originale 488/521
\iffalse Pag. originale 488/521 \fi
\begin{quiz}
Per quali valori di $a\in\mathbb R$ la funzione $f(x)=\log|x-5|$ soddisfa le condizioni del teorema degli zeri, nell\'intervallo $[a,7]$?
\begin{enumerate}
\item $a<6$
\item $a\in[5,6)$
\item $a\in(5,6]$
\item $a\in[5,6]$
\item $a\in(5,6)$
\end{enumerate}
\end{quiz}
% SOLUZIONE: E
% SEZIONE_AM1: 4.1
% SPIEGAZIONE: La funzione deve essere continua su $[a,7]$: occorre $a>5$. Inoltre $f(7)=\log2>0$ e $f(a)<0$ richiede $|a-5|<1$, cio\`e $a<6$. Il caso $a=6$ ha gi\`a uno zero all\'estremo, ma non assicura uno zero interno tramite il criterio dei segni opposti.

% Pag. originale 488/521
\iffalse Pag. originale 488/521 \fi
\begin{quiz}
Il modulo del numero complesso $w=\dfrac{1+i}{3+4i}$ \`e
\begin{enumerate}
\item $\dfrac5{12}$
\item $\dfrac{\sqrt2}{5}$
\item $\sqrt{\dfrac7{12}}$
\item $\sqrt{\dfrac25}$
\item $\dfrac2{25}$
\end{enumerate}
\end{quiz}
% SOLUZIONE: B
% SEZIONE_AM1: 1.3
% SPIEGAZIONE: $|w|=|1+i|/|3+4i|=\sqrt2/5$.

% Pag. originale 489/521
\iffalse Pag. originale 489/521 \fi
\begin{quiz}
$\displaystyle\lim_{x\to+\infty}\frac{\log(x^{-1/2})}{2-\cos^2x}=$
\begin{enumerate}
\item $+\infty$
\item $-\infty$
\item $0$
\item $\nexists$
\item $0^-$
\end{enumerate}
\end{quiz}
% SOLUZIONE: B
% SEZIONE_AM1: 3.4
% SPIEGAZIONE: $\log(x^{-1/2})=-\frac12\log x\to-\infty$, mentre $1\le2-\cos^2x\le2$. Il quoziente tende a $-\infty$.

% Pag. originale 489/521
\iffalse Pag. originale 489/521 \fi
\begin{quiz}
L\'integrale improprio $\displaystyle\int_1^{+\infty}\log\left(1+\frac1{x^4}\right)\,dx=$
\begin{enumerate}
\item converge a zero
\item diverge positivamente
\item converge ad un numero positivo
\item \`e indeterminato
\item $\displaystyle\int_1^{+\infty}\log\left(1+\frac1{x^4}\right)\,dx\sim\int_1^{+\infty}\frac1{x^4}\,dx$
\end{enumerate}
\end{quiz}
% SOLUZIONE: C
% SEZIONE_AM1: 5.5
% SPIEGAZIONE: Per $x\to\infty$, $\log(1+x^{-4})\sim x^{-4}$: l\'integrale converge e, essendo l\'integranda positiva, ha valore positivo. L\'alternativa E usa impropriamente $\sim$ tra valori di integrali definiti, che non sono necessariamente equivalenti.

% Pag. originale 489/521
\iffalse Pag. originale 489/521 \fi
\begin{quiz}
Sia $f(x)=(x-10)^6\cdot\sin x$. Allora
\begin{enumerate}
\item $\exists$ un solo $c\in(0,10)$ tale che $f\prime(c)=0$
\item $\exists$ un solo $c\in(\pi,10)$ tale che $f\prime(c)=0$
\item $x=10$ \`e punto di flesso a tangente orizzontale
\item $x=0$ \`e punto critico
\item $\exists c\in(0,\pi)$ tale che $f\prime(c)=0$
\end{enumerate}
\end{quiz}
% SOLUZIONE: E
% SEZIONE_AM1: 4.3
% SPIEGAZIONE: $f(0)=f(\pi)=0$ e $f$ \`e continua su $[0,\pi]$ e derivabile all\'interno. Per Rolle esiste $c\in(0,\pi)$ con $f\prime(c)=0$.

% Pag. originale 489/521
\iffalse Pag. originale 489/521 \fi
\begin{quiz}
La parte principale, per $x\to0^+$, della funzione $f(x)=\sqrt{x^2+x^3}-x-x^3$ \`e
\begin{enumerate}
\item $2x$
\item $x$
\item $x^2$
\item $\dfrac12x^2$
\item $-x^3$
\end{enumerate}
\end{quiz}
% SOLUZIONE: D
% SEZIONE_AM1: 3.3
% SPIEGAZIONE: Per $x>0$, $\sqrt{x^2+x^3}=x\sqrt{1+x}=x+\frac12x^2+O(x^3)$. Sottraendo $x+x^3$ rimane $\frac12x^2+O(x^3)$.

% Pag. originale 489/521
\iffalse Pag. originale 489/521 \fi
\begin{quiz}
Sia $f(x)=2+3x+\arctan x$. Allora $(f^{-1})\prime(2)=$
\begin{enumerate}
\item $\dfrac{16}{5}$
\item $4$
\item $\dfrac14$
\item $\dfrac12+\dfrac1{3\cdot2}+\dfrac1{\arctan2}$
\item nessuna delle altre risposte \`e esatta
\end{enumerate}
\end{quiz}
% SOLUZIONE: C
% SEZIONE_AM1: 4.2
% SPIEGAZIONE: $f(0)=2$ e $f\prime(0)=3+1=4$. Per il teorema della derivata dell\'inversa, $(f^{-1})\prime(2)=1/4$.

% Pag. originale 489/521
\iffalse Pag. originale 489/521 \fi
\begin{quiz}
Sia $f$ una funzione $\in C^{(n)}(\mathbb R)$ e sia $x_0$ un punto di massimo locale della funzione $f$. Allora \`e vero che
\begin{enumerate}
\item $\exists\delta>0:\ \forall x\in\{|x-x_0|<\delta\}\Rightarrow f(x)<f(x_0)$
\item $f(x)>f(x_0)$, $\forall x\in\mathbb R$
\item $\exists\delta>0:\ \forall x\in\{0<|x-x_0|<\delta\}\Rightarrow f(x)\le f(x_0)$
\item $f(x)\le f(x_0)$, $\forall x\in\mathbb R$
\item $\exists\delta>0:\ \forall x\in\{|x-x_0|<\delta\}\Rightarrow f(x)\le f(x_0)$
\end{enumerate}
\end{quiz}
% SOLUZIONE: E
% SEZIONE_AM1: 4.3
% SPIEGAZIONE: La definizione di massimo locale \`e esattamente l\'alternativa E. Anche C \`e vera, perch\`e restringe la stessa propriet\`a ai soli $x\ne x_0$: il testo presenta due alternative corrette.
Risposta completata
